% GENERATED, edit the French unit or the translation table
\mxunit{terminale-algorithmique}{Algorithmics · P1 to P6}
These programs are examples from the lessons. To understand them, specify the inputs, the output and the hypotheses, then run a few steps by hand. The answers to the activities and homework remain to be written up separately.
\section{P1 · Terms, sum and threshold}
We assume $n$ is a non-negative integer and the safety cap is a positive integer. The first function returns $u_n$ and $\sum_{k=0}^n u_k$ for $u_{k+1}=au_k+b$. The second looks for a first upward crossing; its safety cap only signals that it was not found within the fixed number of steps.
\begin{lstlisting}
def terme_somme(n, u0, a, b):
    u, total = u0, u0
    for k in range(n):
        u = a*u + b
        total += u
    return u, total

def seuil_croissant(u0, a, b, cible, maximum=100000):
    u, n = u0, 0
    while u <= cible and n < maximum:
        u = a*u + b
        n += 1
    if u <= cible:
        return None
    return n
\end{lstlisting}
\section{P2 · Bisection}
The calculations below use floating-point numbers; the mathematical guarantee assumes the signs are correctly determined. Hypotheses: $f$ continuous on $[a,b]$ and $f(a)f(b)\le0$. To identify a unique zero, prove uniqueness separately. The function returns an interval of length at most $\varepsilon$, as long as machine precision still allows the interval to be split.
\begin{lstlisting}
def dichotomie(f, a, b, epsilon):
    assert a < b and epsilon > 0
    fa, fb = f(a), f(b)
    assert fa*fb <= 0
    if fa == 0: return a, a
    if fb == 0: return b, b
    while b-a > epsilon:
        m = (a+b)/2
        if m == a or m == b:
            raise ArithmeticError("precision machine atteinte")
        fm = f(m)
        if fm == 0: return m, m
        if fa*fm < 0:
            b, fb = m, fm
        else:
            a, fa = m, fm
    return a, b
\end{lstlisting}
\mxnewpage\section{P3 · Euler}
With $N$ a strictly positive integer and $t_0<\text{end}$, for $y'=F(t,y)$, $Y_{k+1}=Y_k+hF(t_k,Y_k)$. The program returns all the approximations; it does not prove their accuracy.
\begin{lstlisting}
def euler(F, t0, y0, fin, N):
    assert N > 0 and fin > t0
    h = (fin-t0)/N
    y, valeurs = y0, [(t0, y0)]
    for k in range(N):
        t = t0+k*h
        y = y+h*F(t, y)
        valeurs.append((t0+(k+1)*h, y))
    return valeurs
\end{lstlisting}
\section{P4 · Rectangles}
For $N$ a strictly positive integer and $f$ continuous and increasing on $[a,b]$, the result gives a lower bound and an upper bound for the integral. Without monotonicity, these are only two approximations to be examined.
\begin{lstlisting}
def rectangles(f, a, b, N):
    assert N > 0 and a < b
    h = (b-a)/N
    gauche = h*sum(f(a+k*h) for k in range(N))
    droite = h*sum(f(a+k*h) for k in range(1, N+1))
    return gauche, droite
\end{lstlisting}
\mxnewpage\section{P5 · Binomial distribution and simulation}
We assume $n$ is a non-negative integer, $0\le p\le1$ and $k$ an integer. The direct calculation below suits the modest sizes used in the exercises; for very large parameters, specialised numerical libraries are preferable.
\begin{lstlisting}
from math import comb
from random import Random

def probabilite(n, p, k):
    if k < 0 or k > n: return 0.0
    return comb(n, k)*p**k*(1-p)**(n-k)

def cumul(n, p, k):
    return sum(probabilite(n, p, j)
               for j in range(min(k, n)+1))

def simuler_binomiale(n, p, generateur):
    return sum(generateur.random() < p for _ in range(n))

rng = Random(2026)
\end{lstlisting}
\mxneed{16}\section{P6 · Fluctuations of a frequency}
We assume $N$ is a strictly positive integer and $0\le p\le1$. The seed makes it possible to reproduce the same example. Repeat with other seeds; the law of large numbers does not say that the observed deviation decreases at every step.
\begin{lstlisting}
from random import Random

def frequences(N, p, graine=2026):
    rng = Random(graine)
    succes, resultat = 0, []
    for n in range(1, N+1):
        succes += rng.random() < p
        resultat.append(succes/n)
    return resultat
\end{lstlisting}
